\documentclass{amsbook}

\newtheorem{theorem}{Theorem}[chapter]
\newtheorem{lemma}[theorem]{Lemma}

\theoremstyle{definition}
\newtheorem{definition}[theorem]{Definition}
\newtheorem{example}[theorem]{Example}
\newtheorem{xca}[theorem]{Exercise}

\theoremstyle{remark}
\newtheorem{remark}[theorem]{Remark}

\numberwithin{section}{chapter}
\numberwithin{equation}{chapter}

\begin{document}

% \begin{equation*}
% \begin{split}
% \end{split}
% \end{equation*}

\frontmatter

\title{Minqi's Solutions to the Problems of the National Postgraduate Entrance Examination of China}

\author{Minqi Pan\footnote{Website: http://www.minqi-pan.com}}

\date{\today}

\maketitle

\setcounter{page}{4}

\tableofcontents

\mainmatter

\chapter{Mathematics II, 2018}
\section{Problem 1}
Discover a possible pair of $a$ and $b$ that makes \[
\lim_{x\to 0}\left(e^x+ax^2+bx\right)^{\frac{1}{x^2}}=1.
\]
\begin{proof}
We claim that\[
a=-\frac{1}{2},b=-1.
\] is such a possible pair.

By the definition of the exponential function \[
e^x=\sum_{n=0}^\infty \frac{x^n}{n!},
\] we have\[
e^x-\frac{1}{2}x^2-x=1+\sum_{n=3}^\infty \frac{x^n}{n!}.
\]

Define \[
g(x)=\sum_{n=3}^\infty \frac{x^n}{n!}.
\]

Then \[
\left(e^x-\frac{1}{2}x^2-x\right)^{\frac{1}{x^2}}=\left[\left(1+g(x)\right)^\frac{1}{g(x)}\right]^\frac{g(x)}{x^2}.
\]

By the definition of the number $e$\[
e=\lim_{x\to 0}(1+x)^\frac{1}{x}
\] and the fact that \[
\lim_{x\to 0}g(x)=0,
\] we have \[
\lim_{x\to 0}\left(1+g(x)\right)^\frac{1}{g(x)}=e.
\]

Finally, since \[
\lim_{x\to 0}\frac{g(x)}{x^2}=0,
\] we conclude that \[
\lim_{x\to 0}\left[\left(1+g(x)\right)^\frac{1}{g(x)}\right]^\frac{g(x)}{x^2}=e^0=1.
\]
\end{proof}

\end{document}
